\documentclass[10pt]{article}
\usepackage[margin=1in]{geometry}
\usepackage{amsmath,amssymb,booktabs}
\usepackage[hidelinks]{hyperref}
\title{Understanding a Simple Research Problem}
\author{Alex Researcher \\ Department of Computer Science \\ Example University}
\date{}
\begin{document}
\maketitle
\begin{abstract}
A useful paper starts with a clear question. This editable example introduces a
small estimation problem, explains its assumptions, and reports a reproducible
comparison. Replace the example with your own research.
\end{abstract}
\section{Introduction}
Estimating a signal from noisy measurements connects probability, optimization,
and practical experiments. We study a simple setting in which every modeling
choice can be stated explicitly. The classical treatment of information by
Shannon~\cite{shannon1948} provides useful background.
\section{A simple model}
Suppose the observations are $y_i=\theta+\varepsilon_i$, where the noise terms
are independent, have mean zero, and have variance $\sigma^2$. The sample mean is
\begin{equation}
  \widehat\theta=\frac{1}{n}\sum_{i=1}^n y_i,
  \qquad \mathbb{E}\!\left[(\widehat\theta-\theta)^2\right]=\frac{\sigma^2}{n}.
  \label{eq:error}
\end{equation}
Equation~\eqref{eq:error} predicts how additional observations improve precision.
\section{Results}
Table~\ref{tab:results} shows theoretical values for unit noise variance.
\begin{table}[h]
\centering
\caption{Illustrative theoretical error, not experimental measurements.}
\label{tab:results}
\begin{tabular}{lrr}\toprule
Observations & Mean squared error & Standard error \\\midrule
10 & 0.100 & 0.316 \\
50 & 0.020 & 0.141 \\
100 & 0.010 & 0.100 \\\bottomrule
\end{tabular}
\end{table}
\section{Limitations and next steps}
The independence and constant-variance assumptions may fail in practice.
Check them before applying the estimator to a new dataset.
\bibliographystyle{plain}
\bibliography{references}
\end{document}
